%!TEX root = ../../thesis.tex

\section{Input-Level Action Spaces}\label{sec:results-input-actions}

As outlined in~\cref{sec:targets-input-actions}, we use the input action target
\emph{sequence} as a feasibility test and do not go into further detail.
The results are presented in\todo{cref to figures}.

It is

% Answers RQ1 (\cref{sec:introduction-rqs}). Formulation: \cref{sec:input-actions}.

% TODO: One target and one claim, so keep this section short and let its length signal
% its weight relative to \cref{sec:results-corpus-actions}. Sequence
% (\cref{sec:target-sequence}) is a feasibility test for the implementation, not
% evidence about fuzzing --- say so, and do not let the section drift into fuzzing
% language.
%
% What to report:
%   - solve rate over the seeds, and executions to solve conditional on success. The
%     experiment stops once the target is solved (\cref{sec:methodology-setup}), so
%     there is no parameter study to report and the section must not imply one,
%   - depth curves against the coverage-guided climber and the oracle
%     (\cref{sec:methodology-baselines}). Prefix depth is directly observable here,
%     which is the whole reason the target was built,
%   - the mechanism figures: does the model rank the correct action first at each
%     depth, and hold it? "It works" and "it works for the reason claimed" are
%     separate claims, and the probe suite of \cref{sec:impl-probes} distinguishes
%     them.
%
% Name the baseline and the currency for every claim
% (\cref{sec:methodology-currencies}). Report failed and collapsed runs here rather
% than deferring them all to \cref{chap:failure-modes}; that chapter explains them,
% this one counts them.
%
% Close by handing over to \cref{sec:results-corpus-actions}: the input formulation is
% carried no further, and every fuzzing claim in the thesis is made on the corpus
% targets.
